
% Autoría compartida: Mario coordina la discusión.

\subsection{Comparación de latencia, dispersión y rendimiento}

Los resultados obtenidos permiten identificar diferencias
entre ambas implementaciones bajo las condiciones del
experimento. En las mediciones secuenciales sobre localhost,
TCP presentó menores tiempos de respuesta que gRPC con los
dos tamaños de mensaje evaluados.

Con cargas de 100 bytes, la latencia media fue de
0,122\,ms para TCP y 0,547\,ms para gRPC. Esto significa
que el promedio de gRPC fue aproximadamente 4,48 veces
el de TCP. Al utilizar mensajes de 10\,KiB, las medias
aumentaron a 0,197\,ms y 0,564\,ms, respectivamente,
reduciéndose la relación a aproximadamente 2,86 veces.
Estos valores se presentan en la
Tabla~\ref{tab:estadisticas_rtt}.

También se observó un comportamiento diferente al
incrementar el tamaño de la carga útil. La latencia
media de TCP aumentó alrededor de un 61,5\,\%,
mientras que en gRPC el incremento fue cercano
al 3,1\,\%. Aunque esto muestra una menor variación
de la media de gRPC entre los tamaños probados,
no demuestra que tenga mejor escalabilidad en
cualquier situación. Para establecerlo sería
necesario evaluar más tamaños y condiciones de carga.

La dispersión aporta otra perspectiva. Con mensajes
de 100 bytes, gRPC registró una desviación estándar
de 0,231\,ms, frente a 0,028\,ms en TCP. Además,
el percentil 95 fue de 0,819\,ms para gRPC y
0,185\,ms para TCP. Esto indica que la diferencia
no se limitó a los promedios: también hubo una
mayor variación en los tiempos de gRPC durante
esa condición.

Con mensajes de 10\,KiB, las desviaciones estándar
fueron de 0,039\,ms para TCP y 0,110\,ms para gRPC.
Los percentiles 95 se situaron en 0,282\,ms y
0,758\,ms, respectivamente. Por tanto, TCP mantuvo
tiempos más bajos en las medidas de tendencia
central y en el extremo superior de la
distribución observado mediante el P95.

En cuanto al rendimiento, TCP alcanzó una media
de 576,43 mensajes por segundo con 100 bytes,
frente a 442,64 en gRPC. Para 10\,KiB se
registraron 533,24 y 431,98 mensajes por segundo,
respectivamente. En las condiciones evaluadas,
el rendimiento de TCP fue aproximadamente un
30,2\,\% superior con 100 bytes y un 23,4\,\%
superior con 10\,KiB
(Tabla~\ref{tab:rendimiento}).

Estas diferencias deben interpretarse considerando
cómo se implementaron las mediciones. TCP construye
y serializa directamente sus mensajes JSON, mientras
que gRPC utiliza clases generadas a partir de
Protocol Buffers y realiza llamadas remotas sobre
HTTP/2. Por ello, la comparación refleja el
comportamiento conjunto de las implementaciones,
sus bibliotecas y sus operaciones auxiliares,
no únicamente el costo del protocolo de transporte.

Las pruebas posteriores con tres nodos complementan
estos resultados desde el punto de vista funcional.
En las corridas de veinte intercambios, tanto TCP
como gRPC completaron las solicitudes y sus acuses.
Asimismo, las cuatro capturas concurrentes de
300 solicitudes registraron 1200 eventos por
condición, sin errores causales reportados por
el verificador de Lamport.

Sin embargo, estas pruebas utilizaron otro entorno
y tuvieron como propósito comprobar la comunicación
y el orden causal de los eventos. Por ello, no se
emplearon para ampliar las muestras de latencia
ni para comparar mensajes por segundo con el
benchmark secuencial.

\subsection{Facilidad de desarrollo, tipado y manejo de errores}

% Parte gRPC (Mera).
En gRPC, antes de programar hubo que definir el contrato \texttt{messaging.proto} y generar el código con \texttt{grpcio-tools}, que debe regenerarse ante cualquier cambio. A cambio, el servidor solo implementa \texttt{EnviarMensaje} y el cliente lo invoca desde el stub, sin construir ni leer los mensajes, como sí ocurre en TCP con el JSON y el salto de línea. La principal dificultad fue de organización: para usar el reloj compartido de \texttt{src/common/} hubo que agregar la raíz del repositorio a \texttt{sys.path}.

El contrato garantiza los tipos, por lo que el servidor gRPC no necesita las comprobaciones con \texttt{isinstance} del servidor TCP. Sin embargo, no reemplaza toda la validación: un \texttt{timestamp} negativo llega como entero válido y lo rechaza el reloj, y en proto3 un campo omitido toma su valor por defecto \cite{protobufProto3}, de modo que una solicitud sin \texttt{timestamp} llegó con valor 0 sin ningún error. En los fallos, gRPC devuelve códigos estándar con descripción: el cliente recibió \texttt{INVALID\_ARGUMENT} y el servidor siguió atendiendo, mientras que en TCP una solicitud mal formada cierra la conexión sin informar la causa. Como contrapartida, los mensajes binarios no se pueden leer a simple vista al depurar.

Esta comodidad tiene un costo (tabla~\ref{tab:estadisticas_rtt}): con 100 bytes, la latencia media de gRPC fue unas 4,5 veces la de TCP (0,547\,ms frente a 0,122\,ms). Con 10\,KiB, TCP subió un 61\,\% y gRPC solo un 3\,\%, lo que sugiere que el costo adicional de gRPC es principalmente fijo por llamada.

Desde TCP, la comunicación fue directa de implementar: el cliente abre una conexión con el servidor y la reutiliza para enviar varios mensajes. Sin embargo, como TCP transmite un flujo de bytes, tuvimos que definir en el código cómo separar cada solicitud. Para eso usamos objetos JSON codificados en UTF-8 y un salto de línea al final de cada mensaje. El cliente también construye y analiza manualmente el JSON, mientras que el servidor valida los campos recibidos. Esto nos dio control sobre el intercambio, pero exigió mantener el mismo formato y las mismas reglas en ambos extremos.

El tipado queda a cargo de esas validaciones: por ejemplo, se comprueba que el emisor y el mensaje sean textos, que la marca de tiempo sea un entero y que el acuse tenga los campos esperados. En el manejo de errores, el cliente distingue situaciones como un tiempo de espera agotado, una conexión rechazada u otro error de comunicación. Si la solicitud recibida no tiene el formato esperado, el servidor no devuelve un error estructurado para la aplicación y puede cerrar esa conexión. En esta implementación, TCP requirió más trabajo manual para definir el formato, validar los datos y tratar los errores que la interfaz gRPC generada a partir de su contrato.

\subsection{Contraste con la hipótesis inicial}

La hipótesis planteada antes de realizar las mediciones
establecía que la implementación con sockets TCP
presentaría una latencia media menor que gRPC en
localhost, tanto para mensajes de 100 bytes como
para mensajes de 10240 bytes.

Los resultados respaldaron esta expectativa dentro
de las condiciones estudiadas. Con cargas de 100 bytes,
TCP obtuvo una media de 0,122\,ms, inferior a los
0,547\,ms de gRPC. La misma relación se mantuvo
con 10\,KiB, donde TCP registró 0,197\,ms y
gRPC 0,564\,ms.

Estos resultados cumplen el criterio que se había
establecido previamente para considerar respaldada
la hipótesis. Además, las medianas y los percentiles
95 mantuvieron la misma tendencia, por lo que la
diferencia observada no dependió únicamente de
los valores medios.

No obstante, la hipótesis debe aceptarse solo
dentro de su alcance experimental. Las mediciones
se realizaron en una computadora, con comunicación
local y mediante implementaciones específicas
desarrolladas por el equipo. Tampoco se aplicó
una prueba de significación estadística para
generalizar las diferencias a otras poblaciones
de ejecuciones.

En consecuencia, los resultados respaldan la
hipótesis inicial para nuestro experimento,
pero no permiten afirmar que los sockets TCP
ofrezcan siempre menor latencia que gRPC
en cualquier aplicación distribuida.

\subsection{Aplicación posible al Proyecto Fin de Curso}

El Proyecto Fin de Curso de un integrante del equipo corresponde
al Sistema de Gestión de Refuerzo Académico
(SGRA), una aplicación web orientada a gestionar
solicitudes de refuerzo, disponibilidad docente,
sesiones académicas, notificaciones y seguimiento
de los estudiantes.

Su arquitectura contempla una interfaz desarrollada
con Angular, un backend con Spring Boot y una
base de datos PostgreSQL. La comunicación entre
el frontend y el backend se realiza mediante
una API REST sobre HTTP. Por tanto, la práctica
no implica sustituir automáticamente esta
arquitectura por TCP o gRPC.

Si el SGRA evolucionara hacia servicios de
backend independientes, gRPC podría evaluarse
para la comunicación interna entre componentes.
Por ejemplo, un servicio encargado de gestionar
solicitudes de refuerzo podría consultar a
otro servicio los datos académicos necesarios
para validar asignaturas, docentes o
disponibilidad.

En ese escenario, Protocol Buffers permitiría
establecer contratos explícitos para las
solicitudes y respuestas. Además, los códigos
de estado y los límites de tiempo de gRPC
ofrecerían mecanismos definidos para gestionar
errores de comunicación. Esta sería una
alternativa de diseño que requeriría pruebas
propias antes de adoptarse.

Por otro lado, los sockets TCP podrían utilizarse
en un componente interno muy específico que
necesitara intercambiar mensajes simples mediante
un protocolo controlado por el equipo, como un
mecanismo experimental de supervisión de
conectividad. Sin embargo, habría que implementar
la delimitación, validación, autenticación y
gestión de errores de esos mensajes.

Por estas razones, no sería conveniente reemplazar
la API REST del SGRA por sockets TCP personalizados
solo porque estos presentaron menor latencia en
nuestra práctica. Para operaciones como registrar
solicitudes, consultar historiales o gestionar
usuarios, también deben considerarse la seguridad,
el mantenimiento y la integración con el frontend.

En cuanto al teorema CAP, su aplicación depende
de las garantías de consistencia y disponibilidad
de los datos distribuidos, no de escoger TCP
o gRPC. Si en el futuro el SGRA incorporara
réplicas de información académica o de solicitudes,
sería necesario definir cómo actuar ante
interrupciones de comunicación entre ellas.

Por ejemplo, una solicitud de refuerzo no debería
quedar confirmada dos veces por utilizar copias
inconsistentes de los mismos datos. Este problema
requiere decisiones sobre almacenamiento,
transacciones y coordinación, independientemente
del mecanismo empleado para transportar mensajes.

En definitiva, los resultados experimentales
sirven como referencia para analizar alternativas
de comunicación, pero la elección tecnológica
para el SGRA debe partir de sus requisitos
funcionales, de seguridad y de despliegue.

\subsection{Amenazas a la validez y propuestas de mejora}

La principal limitación del benchmark es que las
mediciones se realizaron sobre localhost. Esto
eliminó factores propios de una red entre equipos,
como la congestión, los retardos de transmisión
entre dispositivos y las pérdidas de conectividad.
Por tanto, las latencias obtenidas no representan
necesariamente las que se observarían en un
despliegue real.

Otro aspecto importante es que las condiciones
de medición no incluyeron exactamente las mismas
operaciones internas. En TCP, la serialización
JSON quedó dentro del intervalo medido, mientras
que en gRPC el objeto de solicitud se construyó
antes de iniciar el cronómetro. Además, cada
implementación utiliza bibliotecas y formatos
de mensajes diferentes. Esto limita la posibilidad
de atribuir las diferencias a una sola causa.

El orden de ejecución también pudo influir,
ya que las pruebas TCP se realizaron antes
que las de gRPC. No se midió de manera
instrumental la actividad de otros procesos
ni se controlaron todas las variaciones
de carga del sistema operativo.

Aunque se obtuvieron 500 mediciones por
condición, estas proceden de cinco corridas
consecutivas de cien solicitudes. Los valores
tomados dentro de una misma corrida pueden
compartir condiciones temporales, por lo que
no deben interpretarse automáticamente como
500 repeticiones completamente independientes.

También se utilizaron únicamente dos tamaños
de carga útil, generados mediante caracteres
ASCII repetidos. Esta decisión permitió
controlar la longitud de los mensajes, pero
no representa la diversidad de estructuras
y contenidos que intercambia una aplicación
como el SGRA.

La prueba concurrente de tres nodos se
realizó por separado en Linux, mientras
que el benchmark de latencia se ejecutó
en Windows. Además, las capturas de carga
concurrente utilizaron un inicio coordinado
mediante barreras. Estas diferencias impiden
atribuir directamente a la concurrencia
los cambios que pudieran existir entre
ambos entornos.

Como mejora, se propone repetir el benchmark
en varios equipos conectados mediante una
red física y evaluar un mayor número de
tamaños de mensaje. También sería conveniente
alternar el orden de ejecución de TCP y gRPC
y establecer puntos de inicio y finalización
de la medición equivalentes.

De igual forma, se podrían repetir las pruebas
concurrentes bajo distintas cantidades de
clientes y registrar tanto su rendimiento como
sus errores. Esto permitiría estudiar el
comportamiento ante cargas mayores, sin
confundir la verificación funcional de los
relojes de Lamport con la medición de latencia.

Además, sería útil incorporar pruebas
controladas de desconexión, retrasos y caída
de nodos durante intercambios activos.
Estas pruebas ayudarían a evaluar aspectos
de fiabilidad que no quedan demostrados por
el éxito de las corridas realizadas.

Las mejoras propuestas no forman parte de
los resultados ya obtenidos. Se presentan
como posibilidades para ampliar la evaluación
y comprobar si las tendencias observadas
se mantienen en otros escenarios.
